%!TeX root = main.tex
%!encoding = UTF-8 Unicode

\section{Requirements for Broadband, Low-Intrusion Current Sensing}

Prior to designing the active current shunt, the measurement-driven requirements are defined.
As in the voltage-probe chapter, the requirements are derived from the EMC-relevant frequency range, practical measurement constraints, and expected disturbance mechanisms in the switching environment.

\subsubsection*{Small Signal Response}
The current-sensing path must provide a bandwidth of at least \SI{500}{MHz} to cover the EMC-relevant spectral content of the commutation current.
This bandwidth requirement excludes commercial Rogowski coils for the present application, as their usable bandwidth is typically limited to approximately \SIrange{20}{50}{MHz}~\cite{xin_review_2022}.

\subsubsection*{Output Interface and Signal Routing}
Analogous to the voltage-probe concept, the current probe must provide a \SI{50}{\ohm} output interface.
This is required to route the measured signal into both low-pass and high-pass processing paths with defined impedance conditions.
Consequently, a purely passive shunt implementation is not sufficient for the intended measurement chain.

\subsubsection*{Isolation Requirement for Combined Measurements}
For standalone current measurements, galvanic isolation is not mandatory because the probe can reference the oscilloscope ground.
However, for simultaneous measurements with LISN disturbance voltage, galvanic isolation becomes necessary because the oscilloscope ground is fixed to the aluminum base plate.
In this case, optical isolation of the measurement link is feasible and compatible with the required \SI{50}{\ohm} signal interface.

\subsubsection*{Measurement Intrusion and Shunt Bandwidth}
The inserted sense resistance must be at least one decade below the on-state resistance \mi{R}{DS,on} of the DUT to avoid materially altering the transient damping behavior.
This is essential because the quality factor of the excited resonance is a key EMC-relevant characteristic.
For the considered \SI{750}{V} and \SI{1200}{V} discrete SiC devices in D2PAK, \mi{R}{DS,on} is in the range of \SIrange{20}{50}{m\ohm}~\cite{robert_bosch_gmbh_datasheet_2023,robert_bosch_gmbh_datasheet_2023-1}.
Accordingly, the sense resistance is limited to \SI{2}{m\ohm} or less.

This low resistance directly conflicts with the bandwidth target in conventional shunt realizations, because each shunt exhibits parasitic inductance.
The resulting shunt bandwidth is

\begin{equation}
    \mi{f}{bw,shunt} = \frac{1}{2\pi \tau_\text{shunt}} = \frac{\mi{R}{shunt}}{2\pi\mi{L}{shunt}} 
\end{equation}

In classical shunt design, larger shunt resistance values support higher bandwidth.
For the present EMC application, this trade-off cannot be exploited because the insertion resistance must remain very small.

\subsubsection*{Parasitic Coupling and Interconnect Length}
As for the voltage probe, parasitic coupling into the sensing path must be minimized, and the electrical interconnect length must be kept short enough to limit transmission-line effects.
For the active current shunt, parasitic coupling is particularly critical because the circuit is placed close to the switching node and operates with higher gain than the active voltage probe.

